\section*{第 \ref{ch:b2:diffeq} 章——微分方程}

\begin{solution}{exo:b2:diffeq:1}
第一个矩阵：\omterm{def:b2:reduction:eigen}{特征值} $1, 2$，\omterm{def:b2:reduction:eigen}{特征向量} $(1,0)$ 与 $(1,1)$。
将 $X_0 = (1,0) = 1\cdot(1,0) + 0\cdot(1,1)$ 分解，其解
为
\[
X(t) = \eu^{t}\begin{pmatrix}1\\ 0\end{pmatrix} .
\]
（初始向量本身就是\omterm{def:b2:reduction:eigen}{特征向量}。）

第二个矩阵：$A = 2I + N$，$N = E_{12}$，$N^2 = 0$，$\eu^{tA} =
\eu^{2t}(I + tN)$，故
\[
X(t) = \eu^{2t}\begin{pmatrix} 1 & t\\ 0 & 1\end{pmatrix}
\begin{pmatrix}1\\ 0\end{pmatrix}
= \eu^{2t}\begin{pmatrix}1\\ 0\end{pmatrix}.
\]
\end{solution}

\begin{solution}{exo:b2:diffeq:2}
$y' = \sin(ty)$：$\abs{\sin(ty) - \sin(tz)} \leq \abs t\,\abs{y -
z}$——关于 $y$ 在每一段时间上一致\omterm{def:b2:metric:continuity}{利普希茨}：
在 $\R$ 上有唯一整体解（对每一段时间应用
该定理）。

$y' = y^2$：仅局部\omterm{def:b2:metric:continuity}{利普希茨}，没有整体定理，事实上
$y(0) = 1$ 的解在 $t = 1$ 处爆破。

$y' = \abs y$：$\abs\cdot$ 是 $1$-利普希茨的，整体存在性与
唯一性成立。取 $y(0) = 0$：$y \equiv 0$（由唯一性）。取 $y(0)
= 1$ 时，$y$ 保持为正（不能穿过零解），于是 $y' =
y$，$y = \eu^t$。
\end{solution}

\begin{solution}{exo:b2:diffeq:3}
若某时刻 $y(t_1) = z(t_1)$，则 $y$ 与 $z$ 满足同一个
$t_1$ 处的柯西问题：由唯一性，$y = z$ 在它们的公共
区间上成立——不同的解永不相交。

对 $y' = y(1-y)$，常数 $0$ 与 $1$ 是解。初值
位于 $\intoo{0}{1}$ 中的解永远不能到达 $0$ 或 $1$
（否则将与常数解相碰）：它始终停留在
$\intoo{0}{1}$ 内，因而是整体的（有界，不发生爆破——例如
由\cref{exo:b2:diffeq:9}的判别准则，或因向量场
在被困的带形区域上有界）。
\end{solution}

\begin{solution}{exo:b2:diffeq:4}
$(A - 2I)^2 = \begin{pmatrix}1 & 1\\ -1 & -1\end{pmatrix}^2 = 0$：
以 $D = 2I$，$N = A - 2I$ 作\omterm{thm:b2:reduction:dunford}{邓福德分解}：
\[
\eu^{tA} = \eu^{2t}\,(I + tN)
= \eu^{2t}\begin{pmatrix} 1 + t & t\\ -t & 1 - t\end{pmatrix}.
\]
常数变易法，取 $B(s) = (\eu^{2s}, 0)^{\mathsf T}$：
\[
X(t) = \int_0^t \eu^{(t-s)A}B(s)\,\dd s
= \eu^{2t}\int_0^t \begin{pmatrix} 1 + (t-s)\\ -(t-s)
\end{pmatrix}\dd s
= \eu^{2t}\begin{pmatrix} t + \frac{t^2}{2}\\[2pt] -\frac{t^2}{2}
\end{pmatrix},
\]
利用 $\eu^{(t-s)A}B(s) = \eu^{2(t-s)}(I + (t-s)N)\,(\eu^{2s},
0)^{\mathsf T} = \eu^{2t}\bigl(1 + (t-s),\, -(t-s)\bigr)^{\mathsf
T}$。（验算：$X(0) = 0$；$X' - AX = (\eu^{2t}, 0)^{\mathsf T}$
求导可得。）
\end{solution}

\begin{solution}{exo:b2:diffeq:5}
\emph{级数：} $A = \alpha I + \beta J$，其中 $J =
\begin{pmatrix}0 & -1\\ 1 & 0\end{pmatrix}$，$J^2 = -I$；两个
加项可交换，故 $\eu^{tA} = \eu^{\alpha t}\,\eu^{\beta tJ}$，
且 $\eu^{\beta t J}$ 的级数按偶次与奇次幂分开，
分别合并为 $\cos(\beta t)I + \sin(\beta t)J$：即所述的旋转伸缩
矩阵。

\emph{复数形式：}将 $(x, y) \in \R^2$ 与 $z = x + \iu y$ 等同；
该系统 $X' = AX$ 写作 $z' = (\alpha + \iu\beta)z$，其
解 $z(t) = \eu^{\alpha t}\eu^{\iu\beta t}z_0$ 恰为
螺线：模为 $\eu^{\alpha t}$，辐角以速率
$\beta$ 增加。
\end{solution}

\begin{solution}{exo:b2:diffeq:6}
令 $v(t) = \eu^{-k(t - t_0)}\int_{t_0}^t u$。则
\[
v'(t) = \eu^{-k(t-t_0)}\Bigl(u(t) - k\int_{t_0}^t u\Bigr)
\leq C\,\eu^{-k(t-t_0)},
\]
由假设，从 $t_0$ 到 $t$ 积分（$v(t_0) = 0$）：
$v(t) \leq \frac{C}{k}\bigl(1 - \eu^{-k(t - t_0)}\bigr)$，即
$k\int_{t_0}^t u \leq C\bigl(\eu^{k(t-t_0)} - 1\bigr)$；代回
假设得：$u(t) \leq C\eu^{k(t-t_0)}$。

唯一性与依赖性：积分方程的两个解 $y, z$
满足
\[
\norm{y(t) - z(t)} \leq \norm{y_0 - z_0} +
k\int_{t_0}^{t}\norm{y - z},
\]
由格朗沃尔不等式（取 $C = \norm{y_0 - z_0}$）得指数
型估计；取 $C = 0$ 即得唯一性。
\end{solution}

\begin{solution}{exo:b2:diffeq:7}
作代换 $y = y_1 z$，其中 $y_1 = \frac{\sin t}{t}$：由一般
降阶公式（\cref{prop:b2:diffeq:secondorder}）给出，对 $u
= z'$，
\[
y_1 u' + \Bigl(2y_1' + \frac{2}{t}\,y_1\Bigr)u = 0
\]
（方程化为 $y'' + \frac2t y' + y = 0$）。计算
$2y_1' + \frac2t y_1 = 2\,\frac{t\cos t - \sin t}{t^2} +
\frac{2\sin t}{t^2} = \frac{2\cos t}{t}$：故
\[
\frac{u'}{u} = -\frac{2\cos t}{t}\cdot\frac{t}{\sin t}
= -2\cot t
\quad\Longrightarrow\quad
u = \frac{1}{\sin^2 t}
\quad (\text{相差一个常数}),
\]
从而 $z = -\cot t$，得 $y_2 = y_1 z = -\frac{\cos t}{t}$。
在 $\intoo{0}{\pi}$ 上的通解：
\[
y(t) = \alpha\,\frac{\sin t}{t} + \beta\,\frac{\cos t}{t} .
\]
（它们是零阶球贝塞尔函数。）
\end{solution}

\begin{solution}{exo:b2:diffeq:8}
\omterm{thm:b2:reduction:dunford}{邓福德分解}：$A = D + N$，其中二者可交换，$D$ \omterm{def:b2:reduction:diag}{可对角化}，
\omterm{def:b2:reduction:eigen}{特征值}与原矩阵相同，$N$ 幂零，于是
\[
\eu^{tA} = \eu^{tD}\,\eu^{tN},
\qquad
\eu^{tN} = \sum_{k < n} \frac{t^kN^k}{k!}
\ \text{（矩阵多项式，变量 } t).
\]
设 $-2\alpha = \max_i \Re\lambda_i < 0$。在使
$D$ 对角化的基下，$\vertiii{\eu^{tD}} \leq \eu^{-2\alpha t}$（矩阵元
$\eu^{t\lambda_i}$ 的模为 $\eu^{t\Re\lambda_i}$）；不同基下的\omterm{def:b2:nvs:norm}{范数}
只相差一个常数。因此
\[
\norm{X(t)} \leq \vertiii{\eu^{tA}}\,\norm{X_0}
\leq C'\,\eu^{-2\alpha t}\,(1 + t)^{n-1}\,\norm{X_0}
\leq C\,\eu^{-\alpha t}\norm{X_0} ,
\]
将多项式吸收进一个指数因子，
（$\eu^{-\alpha t}(1+t)^{n-1} \to 0$，故有界）。
\end{solution}

\begin{solution}{exo:b2:diffeq:9}
设 $y$ 为 $\intco{0}{T}$ 上的极大解，$T \leq \infty$，
并设 $T < \infty$。由积分形式得，对 $t < T$，
\[
\norm{y(t)} \leq \norm{y_0} + \int_0^t \bigl(a\norm{y(s)} +
b\bigr)\dd s
\leq \bigl(\norm{y_0} + bT\bigr) + a\int_0^t\norm{y(s)}\,\dd s ,
\]
由格朗沃尔不等式可得，$\norm{y(t)} \leq (\norm{y_0} +
bT)\,\eu^{aT} =: M$，在 $\intco{0}{T}$ 上解落在一个
紧球内。于是 $y' = f(t, y)$ 在 $T$ 附近有界，故 $y$ 在 $T$ 附近
\omterm{def:b2:metric:continuity}{利普希茨}，并且可\omterm{def:b2:metric:continuity}{连续}延拓到 $T$（由柯西
准则）；在 $(T, y(T))$ 处求解柯西问题，可将 $y$
延拓到 $T$ 之外，与极大性矛盾。因此 $T = \infty$：在线性增长下
不存在有限时间逃逸。
\end{solution}

\begin{solution}{exo:b2:diffeq:10}
$r^2 - 3r + 2$ 的特征根为 $1$ 与 $2$，于是
齐次解为 $a\eu^t + b\eu^{2t}$。试取
$\alpha\eu^t$ 行不通，因为 $\eu^t$ 本身已满足
齐次方程（根 $r = 1$ 与右端项“\omterm{pb:b2:diffeq:1}{共振}”）。
取 $y = \alpha t\,\eu^t$，则 $y' = \alpha(1 +
t)\eu^t$，$y'' = \alpha(2 + t)\eu^t$，并且
\[
y'' - 3y' + 2y = \alpha\eu^t\bigl(2 + t - 3 - 3t + 2t\bigr)
= -\alpha\,\eu^t :
\]
$\alpha = -1$，$y_p = -t\,\eu^t$。通解为 $y = a\eu^t
+ b\eu^{2t} - t\eu^t$。由初值条件 $y(0) = y'(0) = 0$：$a + b
= 0$ 且 $a + 2b - 1 = 0$，得 $b = 1$，$a = -1$：
\[
y(t) = \eu^{2t} - (1 + t)\,\eu^{t} .
\]

\end{solution}

\begin{solution}{exo:b2:diffeq:11}
$A = \lambda I + N$，其中 $N = E_{12} + E_{23}$：$N^2 = E_{13}$，
$N^3 = 0$，且 $\lambda I$ 与 $N$ 可交换：
\[
\eu^{tA} = \eu^{\lambda t}\Bigl(I + tN +
\frac{t^2}{2}N^2\Bigr)
= \eu^{\lambda t}\begin{pmatrix}
1 & t & \frac{t^2}{2}\\
0 & 1 & t\\
0 & 0 & 1
\end{pmatrix}.
\]
解为 $X(t) = \eu^{\lambda t}\bigl(X_0 + tNX_0 +
\frac{t^2}2N^2X_0\bigr)$——每个分量都是 $\eu^{\lambda
t}$ 乘以一个次数 $\leq 2$ 的多项式。次数的界为
幂零指数减一：$\eu^{tN}$ 的级数
在 $N^2$ 处截断。
\end{solution}

\begin{solution}{exo:b2:diffeq:12}
\begin{enumerate}
  \item 若 $X(T) = X(0)$，则 $Y(t) = X(t + T)$ 满足 $Y' =
        AY + B(t + T) = AY + B(t)$，且 $Y(0) = X(0)$：
        由唯一性（\cref{thm:b2:diffeq:linear}）得 $Y = X$，
        即 $X$ 是 $T$ 周期的。反之显然。
  \item 在 $\C$ 上三角化：$A = PT'P^{-1}$，其中 $T'$ 为
        上三角，对角线为 $(\lambda_i)$。上三角矩阵的任意幂
        仍是上三角矩阵，对角线为
        $(\lambda_i^k)$，故 $\eu^{TA} = P\eu^{TT'}P^{-1}$
        同一基下仍为上三角，对角线为
        $(\eu^{T\lambda_i})$，即为\omterm{def:b2:reduction:eigen}{特征值}。于是
        $I - \eu^{TA}$ 可逆，当且仅当对每个\omterm{def:b2:reduction:eigen}{特征值}都有 $\eu^{T\lambda} \neq
        1$，亦即 $T\lambda \notin
        2\iu\pi\Z$，这正是所给假设。
  \item 由常数变易公式，$X(T) = \eu^{TA}X(0) +
        \int_0^T\eu^{(T-s)A}B(s)\dd s$，于是 $X(T) = X(0)$ 化为
        \[
        \bigl(I - \eu^{TA}\bigr)X(0) =
        \int_0^{T}\eu^{(T-s)A}B(s)\,\dd s ,
        \]
        在可逆性假设下，$X(0)$ 有唯一解，
        从而恰有一个 $T$ 周期
        解。对调和振子（$\lambda =
        \pm\iu\omega$），被排除的情形为 $\omega T \in
        2\pi\Z$：外力周期是
        固有周期的倍数——即\omterm{pb:b2:diffeq:1}{共振}，正如周末问题所
        定量刻画。
\end{enumerate}
\end{solution}

\begin{solution}{pb:b2:diffeq:1}
\textbf{1.} \omterm{def:b2:reduction:charpoly}{特征多项式}为 $\lambda^2 -
\tau\lambda + \delta$，根为
$\frac{\tau\pm\sqrt\Delta}{2}$。若 $\delta < 0$，则 $\Delta =
\tau^2 - 4\delta > 0$，两实根之积为
$\delta < 0$，两根异号。若 $\delta > 0$ 且 $\Delta \geq
0$，则实根之积 $> 0$，和为 $\tau$，两根均与
$\tau$ 同号。若 $\Delta < 0$，则为共轭复根对 $\alpha \pm
\iu\beta$，其中 $\alpha = \frac\tau2$，$\beta =
\frac{\sqrt{-\Delta}}2$。

\textbf{2.} $X(t) = a\,\eu^{\mu t}v_- + b\,\eu^{\lambda t}v_+$.
$a = 0$（相应地 $b = 0$）的轨道沿不稳定
（相应地为稳定）特征直线运动；其余轨道在两个时间方向上都有 $\norm X \to \infty$，
并且渐近于 $\R v_+$（当 $t \to
+\infty$ 时），渐近于 $\R v_-$（当 $t \to -\infty$ 时）：
这就是鞍点图像。整个 $\R$ 上的有界性迫使 $b = 0$（否则在 $+\infty$ 处爆炸）
且 $a = 0$（在 $-\infty$ 处），只有
原点满足要求。

\textbf{3.} 因 $\mu < \lambda < 0$，两个指数项都衰减：
$X(t) \to 0$。若 $b \neq 0$，提出因子 $\eu^{\lambda t}$：
\[
X(t) = \eu^{\lambda t}\bigl(b\,v_\lambda + a\,\eu^{(\mu -
\lambda)t}v_\mu\bigr),
\qquad \eu^{(\mu-\lambda)t} \to 0 :
\]
$X(t)$ 的方向趋于 $\R v_\lambda$，
\emph{慢}特征方向——除快轴外的所有轨道
都与之相切地趋近（参见本章相图的
右侧子图）。

\textbf{4.} 在使 $A = \begin{pmatrix} \alpha &
-\beta\\ \beta & \alpha\end{pmatrix}$
的基下（见\cref{exo:b2:diffeq:5}；对第二部分的振子系统，
此形式可经显式实基变换得到），
解为 $\eu^{\alpha t}$ 与旋转角为 $\beta
t$ 的旋转之积，即对数螺线，当 $\alpha = \frac\tau2
< 0$ 时收缩，当 $\tau > 0$ 时扩张，而闭曲线（
原坐标下为椭圆）出现在 $\tau = 0$ 时：中心。

\textbf{5.} $\Delta = 0$ 给出二重\omterm{def:b2:reduction:eigen}{特征值} $\lambda =
\frac\tau2$；由凯莱–哈密顿
（\cref{thm:b2:reduction:cayleyhamilton}），$(A - \lambda I)^2 =
0$，故 $N = A - \lambda I$ 幂零，并与
$\lambda I$ 可交换，且 $\eu^{tA} = \eu^{\lambda t}(I + tN)$。若 $N =
0$：$A = \lambda I$，所有射线都是轨道（星形结点）。若 $N
\neq 0$：$X(t) = \eu^{\lambda t}(X_0 + tNX_0)$，且对 $NX_0
\neq 0$，方向收敛于唯一的特征方向
$\operatorname{im}N$：非正常结点。至此完成了
迹–\omterm{def:b2:linalg:det}{行列式}图像。

\textbf{6.} $\tau = -2\zeta\omega$，$\delta = \omega^2 > 0$，
$\Delta = 4\omega^2(\zeta^2 - 1)$。于是，$0 < \zeta < 1$ 时
$\Delta < 0$，$\tau < 0$：稳定焦点；当 $\zeta = 1$ 时：$\Delta =
0$：退化稳定结点；$\zeta > 1$ 时：$\Delta > 0$，$\tau <
0$，$\delta > 0$：稳定结点；$\zeta = 0$ 时：$\tau = 0$，
$\delta > 0$：中心。平面中的一段竖直路径，位于
$\delta = \omega^2$ 处。

\textbf{7.} 根为 $r = -\zeta\omega \pm
\omega\sqrt{\zeta^2-1}$。当 $\zeta < 1$ 时：$r = -\zeta\omega
\pm \iu\omega_d$，$\omega_d = \omega\sqrt{1-\zeta^2}$：
\[
x(t) = \eu^{-\zeta\omega t}\bigl(a\cos\omega_dt +
b\sin\omega_dt\bigr)
= R\,\eu^{-\zeta\omega t}\cos(\omega_dt - \varphi) .
\]
当 $\zeta = 1$ 时：$x = (a + bt)\eu^{-\omega t}$。当 $\zeta >
1$ 时：$x = a\eu^{r_-t} + b\eu^{r_+t}$，两个速率均为负。
欠阻尼情形下，$\abs x$ 相继的极大值出现于
彼此相隔一个伪周期 $\frac{2\pi}{\omega_d}$
（余弦相位相同），它们的比值为
$\eu^{-\zeta\omega\cdot2\pi/\omega_d} =
\eu^{-2\pi\zeta/\sqrt{1-\zeta^2}}$：即对数衰减率，
一个可在示波器上读出的阻尼指标。

\textbf{8.} 有理化，
\[
\abs{\lambda_{\mathrm{slow}}} = \omega\bigl(\zeta -
\sqrt{\zeta^2-1}\bigr)
= \frac{\omega}{\zeta + \sqrt{\zeta^2 - 1}} ,
\]
其分母随 $\zeta \geq 1$ 增大，因此衰减
速率在 $\zeta = 1$ 时最大，且等于 $\omega$。一个
过阻尼的门关闭时不会猛然撞上，但是\emph{缓慢}；
临界阻尼才是工程上的最优选择。

\textbf{9.} $E' = x'x'' + \omega^2xx' =
x'\bigl(-2\zeta\omega x' - \omega^2x\bigr) + \omega^2xx' =
-2\zeta\omega\,x'^2 \leq 0$。若 $x$ 是周期的，且
非常数，则 $E$ 将是周期且不增的，因而
为常数，迫使 $x' \equiv 0$：$x$ 为常数，于是
$\omega^2x = 0$，故 $x \equiv 0$。因此当 $\zeta > 0$ 时，唯一的
周期解为静止解：阻尼消灭了每一个周期。

\textbf{10.} 中心 $\zeta = 0$ 位于 $\tau =
0$ 的直线上，是迹–\omterm{def:b2:linalg:det}{行列式}平面中一个内部为空的
集合：矩阵的任意小扰动
（任何物理阻尼皆然）都会使 $\tau$ 离开零，并把
闭轨道变为螺旋线。无阻尼振子的周期性
因而只是一种刀锋式的现象，而非稳健的
性质。

\textbf{11.} 将 $x_p = \Re(z\eu^{\iu\gamma t})$ 代入
方程，得
\[
\bigl(-\gamma^2 + 2\iu\zeta\omega\gamma +
\omega^2\bigr)z = F
\quad\Longrightarrow\quad
z = \frac{F}{\omega^2 - \gamma^2 + 2\iu\zeta\omega\gamma},
\]
于是 $x_p = \abs z\cos(\gamma t - \varphi)$，其中 $\varphi =
\arg(\omega^2 - \gamma^2 + 2\iu\zeta\omega\gamma)$，即\
$\tan\varphi = \frac{2\zeta\omega\gamma}{\omega^2 -
\gamma^2}$，且所述的 $A(\gamma) = \abs z$。

\textbf{12.} 两解之差满足
齐次方程，当 $\zeta > 0$ 时它衰减到 $0$
（第 7 题）：每个解都等于 $x_p$ 加上一个瞬态项，
该项在无穷远处趋于零。稳态是一个全局吸引子：
初值被遗忘，仅 $A(\gamma)$ 与
相位滞后 $\varphi$ 保持不变。

\textbf{13.} 求 $g(u) = (\omega^2 - u)^2 +
4\zeta^2\omega^2u$ 在 $u = \gamma^2 \geq 0$ 上的最小值。$g'(u) =
-2(\omega^2 - u) + 4\zeta^2\omega^2 = 0$ 的解为 $u = \omega^2(1 -
2\zeta^2)$，内点当且仅当 $\zeta < \frac{1}{\sqrt2}$。此时
\[
g(u_*) = 4\zeta^4\omega^4 + 4\zeta^2\omega^4(1 - 2\zeta^2)
= 4\zeta^2\omega^4(1 - \zeta^2),
\qquad
A(\gamma_*) = \frac{F}{2\zeta\omega^2\sqrt{1 - \zeta^2}} .
\]
与静态响应 $A(0) = \frac{F}{\omega^2}$ 相比：
放大倍数 $\frac{1}{2\zeta\sqrt{1-\zeta^2}} \approx
\frac{1}{2\zeta}$（$\zeta$ 较小时）。弱阻尼
系统在 $\gamma_* \approx \omega$ 附近会把输入放大，
当 $\zeta = 0.005$ 时可达百倍。

\textbf{14.} 所给的 $x$ 满足 $x(0) = x'(0) = 0$，且
\[
x'' + \omega^2 x = \frac{F(\omega^2 -
\gamma^2)\cos\gamma t}{\omega^2 - \gamma^2} = F\cos\gamma t
\]
（$\cos\omega t$ 项相消）。乘积形式
由和差化积公式 $\cos p - \cos q = 2\sin\frac{q+p}{2}\sin\frac{q-p}2$
（取 $p = \gamma t$，$q = \omega t$）推出。当 $\gamma$ 接近
$\omega$ 时，因子 $\sin\frac{(\omega-\gamma)t}2$ 是一个慢变的
包络，调制快速振荡
$\sin\frac{(\omega+\gamma)t}2$：拍，振幅为
$\frac{2F}{\abs{\omega^2-\gamma^2}}$，很大但有界。

\textbf{15.} 对 $x_p = \frac{F}{2\omega}t\sin\omega t$：
\[
x_p'' = \frac{F}{2\omega}\bigl(2\omega\cos\omega t -
\omega^2t\sin\omega t\bigr)
= F\cos\omega t - \omega^2x_p :
\]
为一个解。对固定的 $t$，令 $\gamma \to \omega$，
代入第 14 题的解，取极限：
\[
\frac{2F\sin\frac{(\omega-\gamma)t}2
\sin\frac{(\omega+\gamma)t}{2}}
{(\omega-\gamma)(\omega+\gamma)}
\longrightarrow
\frac{2F\cdot\frac{(\omega-\gamma)t}2\big/(\omega-\gamma)
\cdot\sin\omega t}{2\omega}
= \frac{F\,t\sin\omega t}{2\omega} .
\]
振幅无界地线性增长：\omterm{pb:b2:diffeq:1}{共振}
灾难——这正是士兵过桥时要打乱步伐的原因。

\textbf{16.} 方波含有每个奇数频率
$n = 1, 3, 5, \dots$ 的谐波；由线性性，每一谐波
$n$ 被振子在 $\gamma = n$ 处的响应放大。
当 $\omega = 3$ 时，第三谐波恰好发生\omterm{pb:b2:diffeq:1}{共振}。
工程师忌惮方波（及锯齿波）输入，因为它们
同时激发\emph{全部}奇次谐波：无论结构的固有
频率为何，总有某个谐波在那里等着它。

\textbf{17.} $W' = y_1y_2'' - y_1''y_2 = -qy_1y_2 + qy_1y_2 =
0$：$W$ 为常数（刘维尔公式，伴随矩阵
的迹为零）。$W = 0$ 在某点成立时，$y_2$ 的初值
与 $y_1$ 的初值成比例，从而 $y_2$ 与
$y_1$ 成比例（由唯一性）；$W \neq 0$ 当且仅当线性无关。若 $y(t_0) =
y'(t_0) = 0$，则 $y \equiv 0$（由唯一性）：非零
解的零点都是单零点，而单零点是孤立的
（附近 $y'$ 保号）。

\textbf{18.} 在相邻零点 $a < b$ 之间，$y_1$ 保持
符号不变，不妨设 $y_1 > 0$ 于 $\intoo ab$ 上，则 $y_1'(a) > 0$ 且
$y_1'(b) < 0$（单零点）。将常数 $W =
y_1y_2' - y_1'y_2$ 分别在 $a$ 与 $b$ 处求值：
\[
W = -y_1'(a)\,y_2(a) = -y_1'(b)\,y_2(b) ,
\]
于是 $y_2(a)$ 与 $y_2(b)$ 异号（$W \neq 0$
排除二者为零）：$y_2$ 在 $\intoo ab$ 内有零点
（由介值定理）。该区间内不能有两个零点：
若 $y_2$ 的两个零点夹住 $y_1$ 的一个零点，
由同样论证（交换角色）与相邻性矛盾，故恰有
一个零点——即交错性。

\textbf{19.} 设 $z$ 在 $\intoo ab$ 内无零点；用相反数替换
$y, z$ 后不妨设 $y > 0$ 且 $z > 0$ 于
$\intoo ab$ 上。令 $\varphi = yz' - y'z$，则 $\varphi' = yz'' -
y''z = (q_1 - q_2)\,yz \leq 0$ 于 $\intoo ab$ 上，$\varphi$ 为
不增函数。但 $\varphi(a) = -y'(a)z(a) \leq 0$（因为
$y'(a) > 0$，$z(a) \geq 0$），而 $\varphi(b) = -y'(b)z(b) \geq
0$（因为 $y'(b) < 0$，$z(b) \geq 0$）：
这一不增函数自 $\leq 0$ 变至 $\geq 0$，必恒为零，故
$(q_1 - q_2)yz \equiv 0$ 于 $\intoo ab$ 上。若 $q_1 < q_2$
在 $\intoo ab$ 的某处成立，这是荒谬的（那里 $y, z > 0$）：
$z$ 必在内部严格为零。一般地（$q_1 \leq q_2$），
或 $z$ 在 $\intoo ab$ 内取零，或 $\varphi \equiv 0$
蕴含 $z$ 与 $y$ 成比例，而后者在 $a$ 与 $b$ 处为零：
总之，$z$ 在 $\intcc ab$ 中必有零点。

\textbf{20.} 上界：比较 $y$（系数 $q \geq
m^2$）与 $u(t) = \sin(m(t - a))$（系数 $m^2 \leq q$，
故 $y$ 在第 19 题中扮演 $z$ 的角色）：若 $y$
在 $\intoc{a}{a + \pi/m}$ 内无零点，则 $a$ 与 $a +
\frac\pi m$ 为 $u$ 的相邻零点，二者之间 $y \neq 0$，
这与第 19 题矛盾：
$y$ 的相邻零点距离 $\leq \frac\pi m$。下界：若两个
相邻零点 $a < b$ 为 $y$ 的零点，且 $b - a < \frac\pi M$，
则 $z(t) = \sin(M(t-a))$（系数 $M^2 \geq q$）必
在 $\intcc ab \subset \intoo{a}{a + \pi/M}
\cup\{a\}$ 中取零，该集合中唯一的零点是 $a$ 本身——
但对 $\intoo{a}{b}$ 应用第 19 题，并计及端点处的严格性，
得 $\intcc ab$ 中有零点，而 $z > 0$ 于
$\intoc ab$ 上：矛盾。因此 $\frac\pi M \leq b - a \leq
\frac\pi m$；当 $q \equiv \omega^2$ 时，两个界
合并为谐振子的
精确间距 $\frac\pi\omega$。

\textbf{21.} 令 $u = ty$，则 $u'' = ty'' + 2y'$，故 $ty'' + 2y'
+ ty = u'' + u = 0$：$u = A\sin t + B\cos t = R\sin(t +
\varphi)$，又 $y = \frac{u}{t}$ 给出 $\frac{\sin t}t$ 与
$\frac{\cos t}t$（见\cref{exo:b2:diffeq:7}），无须
降阶。任一非零解的零点与
$R\sin(t + \varphi)$ 的零点相同，间距恰为 $\pi$——
施图姆思想的体现：系数 $q \equiv 1$ 决定
零点，无论是否有公式。

\textbf{22.} 设 $x(t) =
\frac1\omega\int_0^t\sin(\omega(t-s))F(s)\dd s$。则 $x(0) =
0$;
\[
x'(t) = \frac1\omega\sin(0)F(t) +
\int_0^t\cos(\omega(t-s))F(s)\dd s
= \int_0^t\cos(\omega(t-s))F(s)\dd s ,
\]
故 $x'(0) = 0$；且 $x''(t) = F(t) -
\omega\int_0^t\sin(\omega(t-s))F(s)\dd s = F(t) - \omega^2x(t)$
（对变上限\omterm{thm:b2:integration:continuity}{含参积分}求导，
如积分一章所述）。取 $F(s) = F\cos\omega s$，
由积化和差公式得
\[
\int_0^t\sin(\omega(t-s))\cos(\omega s)\dd s
= \frac12\int_0^t\bigl(\sin\omega t + \sin(\omega t -
2\omega s)\bigr)\dd s
= \frac{t}{2}\sin\omega t ,
\]
（第二项积分为零），故 $x =
\frac{F}{2\omega}t\sin\omega t$：第 15 题再次出现，
来自杜哈梅原理。

\textbf{23.} 写成方程组形式 $X' = AX + (0, F(t))^{\mathsf T}$，
其中 $\operatorname{Sp}A$ 的谱实部皆为负（$\zeta >
0$）：由常数变易法及\cref{exo:b2:diffeq:8}
（$\vertiii{\eu^{tA}} \leq C\eu^{-\alpha t}$）给出
\[
\norm{X(t)} \leq C\eu^{-\alpha t}\norm{X_0}
+ \int_0^t C\eu^{-\alpha(t-s)}\norm{F}_\infty\dd s
\leq C\norm{X_0} + \frac{C\norm F_\infty}{\alpha} :
\]
有界输入，有界输出——在暂态过后，对初始
数据一致成立。

\textbf{24.} 当 $\zeta = 0$ 且 $\gamma \neq \omega$ 时，
第 14 题的解有界，加上任一齐次解
（有界：中心的轨道为圆）后，仍
保持有界；当 $\gamma = \omega$ 时，第 15 题的解线性增长。故
对无阻尼振子，周期性外力下的有界性
恰在一个频率处失效——\omterm{pb:b2:diffeq:1}{共振}——
而第 23 题表明，任何正阻尼都能恢复
对\emph{所有}有界输入的有界性。

\textbf{25.} (i) 迹--行列式平面对一切自治平面线性动力系统分类，
而第二部分的振子沿其中一条竖直线运动；
外力使系统离开该平面
（非自治），由杜哈梅原理接手。(ii) $\gamma_*$ 是
系统偏好的频率，$A(\gamma_*)$ 是激发它的代价，
$\frac{1}{2\zeta}$ 是放大因子
——即工程师所称的
品质因数，亦即\omterm{pb:b2:diffeq:1}{共振}的尖锐程度。(iii) 施图姆定理仅凭 $q$ 的符号与大小
即可判定振动情形：它们适用于解没有初等公式的
方程，例如贝塞尔方程与薛定谔方程。
(iv) \omterm{def:b2:diffeq:wronskian}{朗斯基行列式}为常数（第 17 题，经刘维尔公式）与
有界输入稳定性（第 23 题）被不加说明地反复使用，
每当本书遇到变系数方程或
扰动系统时。
\end{solution}
